\documentclass[UTF8,a4paper,12pt]{ctexart}
\usepackage{xeCJK}
\setCJKmainfont{Noto Serif CJK SC}
% ===================== 基本版式 =====================
\usepackage{geometry}
\usepackage{graphicx}
\usepackage{float}
\usepackage{booktabs}
\usepackage{array}
\usepackage{tabularx}
\usepackage{longtable}
\usepackage{enumitem}
\usepackage{hyperref}
\usepackage{xcolor}
\usepackage{amsmath}
\usepackage{tikz}
\usetikzlibrary{arrows.meta, positioning, shapes.geometric}

\geometry{left=2.5cm,right=2.5cm,top=2.5cm,bottom=2.5cm}
\linespread{1.5}
\setlength{\parindent}{2em}
\setlength{\parskip}{0pt}
\setlist[itemize]{noitemsep, topsep=2pt}
\setlist[enumerate]{noitemsep, topsep=2pt}

\ctexset{
  section={format=\zihao{3}\heiti\bfseries, aftername=\quad},
  subsection={format=\zihao{4}\heiti\bfseries, aftername=\quad},
  subsubsection={format=\zihao{-4}\heiti\bfseries, aftername=\quad}
}

\hypersetup{colorlinks=true, linkcolor=black, urlcolor=blue}
\newcommand{\todo}[1]{\textcolor{red}{\textbf{【待补充：#1】}}}
\newcommand{\code}[1]{\texttt{#1}}

\begin{document}

% ===================== 封面 =====================
\begin{titlepage}
    \centering
    \vspace*{1.8cm}
    {\zihao{2}\heiti 第十届全国大学生集成电路创新创业大赛\\[0.4em] ``竞业达''企业命题初赛}\par
    \vspace{1.8cm}
    {\zihao{1}\heiti RISC-V CPU设计报告}\par
    \vspace{0.8cm}
    {\zihao{3}\heiti 基于RV32I指令集的五级流水线CPU设计}\par
    \vfill
    \begin{tabular}{rl}
        院\quad 校： & {上海科技大学} \\[0.6em]
        作品名称： & 基于RV32I指令集的五级流水线CPU设计 \\[0.6em]
        团队名称： & {决战内存条队} \\[0.6em]
        姓\quad 名： & {李建昊、施云翀、袁源} \\[0.6em]
        指导教师： & {刘闯} \\[0.6em]
        时\quad 间： & {2026年5月7日} 
    \end{tabular}
    \vfill
\end{titlepage}

\tableofcontents
\newpage

% ===================== 快速预览 =====================
\section*{作品快速预览简介}
\addcontentsline{toc}{section}{作品快速预览简介}

本作品面向``竞业达''企业命题初赛中基于RISC-V指令集的CPU设计任务，实现了一款基于RV32I基本整数指令集的32位五级流水线CPU。CPU顶层模块为\code{myCPU}，采用SystemVerilog描述，能够接入企业提供的工程模板，通过\code{irom\_addr/irom\_data}访问指令存储器，通过\code{perip\_addr/perip\_wen/perip\_mask/perip\_wdata/perip\_rdata}访问DRAM与外设。

作品关键技术包括：五级流水线数据通路、RV32I指令译码、立即数生成、寄存器堆双读单写、ALU算术逻辑运算、访存数据扩展、EX/MEM/WB阶段数据前递、load-use冒险暂停以及分支/跳转控制冒险冲刷。当前设计以初赛要求的RV32I 37条基础指令为实现目标，\code{fence}、\code{ecall}和\code{ebreak}暂不作为初赛核心测试指令实现。

\begin{table}[H]
\centering
\caption{作品快速预览表}
\begin{tabularx}{\textwidth}{p{4cm}X}
\toprule
项目 & 内容 \\
\midrule
CPU类型 & 32位RISC-V CPU \\
指令集目标 & RV32I基本整数指令集，初赛重点支持37条基础指令 \\
流水线结构 & IF、ID、EX、MEM、WB五级流水线 \\
核心模块 & \code{PC}、\code{NPC}、\code{Control}、\code{IMMGEN}、\code{RegFile}、\code{ALU}、\code{Mask}、流水线寄存器、前递单元、冒险检测单元 \\
冒险处理 & 数据前递、load-use暂停、分支/跳转冲刷 \\
外设接口 & IROM接口与DRAM/外设统一访问接口 \\
Vivado版本 & Vivado 2023.2
开发板型号 & Kintex-7
补充内容 & \todo{整体架构图截图、综合实现结果、bit文件信息} \\
\bottomrule
\end{tabularx}
\end{table}

\newpage

\section{项目概述}

\subsection{项目背景}

RISC-V是一种开放的指令集架构，具有结构简洁、模块化程度高、扩展性强等特点，适合用于处理器设计教学、嵌入式系统开发和FPGA原型验证。本项目围绕``竞业达''企业命题初赛要求，基于RV32I基本整数指令集设计并实现一款32位CPU。通过该设计，可以完整训练从指令编码、数据通路、控制通路、流水线寄存器、冒险处理到FPGA接口适配的处理器设计流程。

本设计采用五级流水线结构，将指令执行过程划分为取指、译码、执行、访存和写回五个阶段。与单周期CPU相比，五级流水线能够在不同阶段并行处理多条指令，提高指令吞吐率；同时，流水线结构也引入了数据冒险和控制冒险问题，因此本设计进一步实现了前递单元和冒险检测单元，以保证程序执行结果的正确性。

\subsection{设计目标}

本项目的设计目标如下：

\begin{enumerate}
    \item 实现一款基于RV32I基本整数指令集的32位CPU；
    \item 采用五级流水线结构，实现IF、ID、EX、MEM、WB五个阶段；
    \item 支持初赛要求的37条RV32I基础指令，不将\code{fence}、\code{ecall}和\code{ebreak}作为初赛核心测试对象；
    \item 完成PC更新、指令译码、立即数生成、寄存器读写、ALU运算、访存和写回等功能；
    \item 实现数据前递、load-use暂停和控制冒险冲刷机制；
    \item 适配企业工程模板提供的\code{myCPU}接口；
    \item 通过外设统一接口完成对DRAM、数码管、LED等地址空间的访问；
    \item 为线上测评、上板验证、bit文件复核和演示视频提交提供硬件基础。
\end{enumerate}

\subsection{设计平台说明}

\begin{table}[H]
\centering
\caption{设计平台说明}
\begin{tabularx}{\textwidth}{p{4cm}X}
\toprule
项目 & 内容 \\
\midrule
硬件平台 & {竞业达数字孪生平台} \\
FPGA器件 & {Xilinx Kintex-7} \\
开发工具 & {Vivado 2023.2} \\
硬件描述语言 & SystemVerilog \\
仿真工具 & {Vivado Simulation、Verilator} \\
CPU时钟 & 企业模板默认提供50MHz外设时钟，可通过PLL调整；本设计使用CPU时钟频率为{130MHz} \\
工程顶层 & 企业模板顶层 + 自研\code{myCPU}核心 \\
生成文件 & {src0.bit, src1.bit, src2.bit} \\
\bottomrule
\end{tabularx}
\end{table}

\section{RISC-V CPU架构设计}

\subsection{RV32I指令集的支持情况}

本CPU面向RV32I基本整数指令集进行设计。初赛测试要求中，RV32I共40条基础指令，排除\code{fence}、\code{ecall}、\code{ebreak}三条后，最低要求实现37条指令。本设计按照该目标对R型、I型、Load、Store、Branch、U型和J型指令进行支持。

\begin{longtable}{p{3cm}p{8cm}p{3cm}}
\caption{RV32I指令支持情况}\\
\toprule
指令类型 & 指令 & 当前设计状态 \\
\midrule
\endfirsthead
\toprule
指令类型 & 指令 & 当前设计状态 \\
\midrule
\endhead
R型运算指令 & \code{add, sub, sll, slt, sltu, xor, srl, sra, or, and} & 已设计 \\
I型运算指令 & \code{addi, slti, sltiu, xori, ori, andi, slli, srli, srai} & 已设计 \\
Load指令 & \code{lb, lh, lw, lbu, lhu} & 已设计 \\
Store指令 & \code{sb, sh, sw} & 已设计 \\
Branch指令 & \code{beq, bne, blt, bge, bltu, bgeu} & 已设计 \\
U型指令 & \code{lui, auipc} & 已设计 \\
J型指令 & \code{jal, jalr} & 已设计  \\
\bottomrule
\end{longtable}

\subsection{RV32I CPU的整体架构图}

本设计采用五级流水线结构。整体架构由取指阶段、译码阶段、执行阶段、访存阶段和写回阶段组成，各阶段之间通过流水线寄存器传递数据与控制信号。
\begin{figure}[H]
\centering
\resizebox{0.96\textwidth}{!}{
\begin{tikzpicture}[
    font=\small,
    stage/.style={
        draw,
        rounded corners,
        minimum width=1.65cm,
        minimum height=0.9cm,
        align=center,
        thick
    },
    reg/.style={
        draw,
        minimum width=0.75cm,
        minimum height=0.75cm,
        align=center,
        font=\scriptsize
    },
    unit/.style={
        draw,
        rounded corners,
        minimum width=2.05cm,
        minimum height=0.65cm,
        align=center,
        font=\scriptsize
    },
    hazard/.style={
        draw,
        rounded corners,
        minimum width=3.2cm,
        minimum height=0.65cm,
        align=center,
        font=\scriptsize
    },
    dataarrow/.style={
        -{Stealth[length=2mm]},
        thick
    },
    ctrlarrow/.style={
        -{Stealth[length=2mm]},
        dashed,
        thick,
        rounded corners=3pt
    }
]

% ================= 主流水线阶段 =================
\node[stage] (IF) at (0,0) {IF\\取指};
\node[reg]   (IFID) at (1.75,0) {IF/ID};
\node[stage] (ID) at (3.50,0) {ID\\译码};
\node[reg]   (IDEX) at (5.25,0) {ID/EX};
\node[stage] (EX) at (7.00,0) {EX\\执行};
\node[reg]   (EXMEM) at (8.85,0) {EX/MEM};
\node[stage] (MEM) at (10.75,0) {MEM\\访存};
\node[reg]   (MEMWB) at (12.65,0) {MEM/WB};
\node[stage] (WB) at (14.40,0) {WB\\写回};

% ================= 主数据通路 =================
\draw[dataarrow] (IF.east) -- (IFID.west);
\draw[dataarrow] (IFID.east) -- (ID.west);
\draw[dataarrow] (ID.east) -- (IDEX.west);
\draw[dataarrow] (IDEX.east) -- (EX.west);
\draw[dataarrow] (EX.east) -- (EXMEM.west);
\draw[dataarrow] (EXMEM.east) -- (MEM.west);
\draw[dataarrow] (MEM.east) -- (MEMWB.west);
\draw[dataarrow] (MEMWB.east) -- (WB.west);

% ================= 各阶段内部模块 =================
\node[unit] (ifunit) at (0,-1.45) {PC / NPC\\IROM接口};
\node[unit] (idunit) at (3.50,-1.45) {Control\\IMMGEN / RegFile};
\node[unit] (exunit) at (7.00,-1.45) {ALU\\ALU\_src};
\node[unit] (memunit) at (10.75,-1.45) {DRAM/外设接口\\Mask};
\node[unit] (wbunit) at (14.40,-1.45) {写回选择\\RegFile写入};

% 各阶段模块到对应流水线阶段，短虚线，不跨框
\draw[ctrlarrow] (ifunit.north) -- (IF.south);
\draw[ctrlarrow] (idunit.north) -- (ID.south);
\draw[ctrlarrow] (exunit.north) -- (EX.south);
\draw[ctrlarrow] (memunit.north) -- (MEM.south);
\draw[ctrlarrow] (wbunit.north) -- (WB.south);

% ================= 冒险处理单元 =================
\node[hazard] (hazardunit) at (7.00,-2.85) {Forwarding Unit\\Hazard Detection Unit};

% 冒险处理到EX阶段，短虚线
\draw[ctrlarrow] (hazardunit.north) -- (exunit.south);

% 冒险处理到IF/ID，路线从底部绕行，不压框
\draw[ctrlarrow]
    (hazardunit.west)
    -- (1.75,-2.85)
    -- (IFID.south);

% 冒险处理到ID/EX，路线走模块间隙，不压框
\draw[ctrlarrow]
    (hazardunit.west)
    -- (5.25,-3.25)
    -- (IDEX.south);

% 写回反馈到寄存器堆，从右侧绕到底部，不穿过中间框图
\draw[ctrlarrow]
    (WB.east)
    -- ++(0.45,0)
    |- (3.50,-4.05)
    -- (idunit.south);

\end{tikzpicture}
}
\caption{五级流水线 CPU 整体架构示意图}
\label{fig:cpu_arch}
\end{figure}

\todo{如需使用工程真实框图，可在此处替换为Vivado RTL Analysis或draw.io绘制的整体架构图。}

\subsection{CPU顶层接口设计}

本CPU顶层模块为\code{myCPU}，与企业模板中\code{Core\_cpu}实例连接。接口包括系统时钟复位、IROM指令存储器接口以及DRAM/外设统一接口。

\begin{table}[H]
\centering
\caption{\code{myCPU}顶层接口说明}
\begin{tabularx}{\textwidth}{p{3cm}p{2cm}p{2cm}X}
\toprule
信号名称 & 方向 & 位宽 & 功能说明 \\
\midrule
\code{cpu\_clk} & input & 1 & CPU工作时钟 \\
\code{cpu\_rst} & input & 1 & CPU复位信号 \\
\code{irom\_addr} & output & 32 & 指令存储器访问地址，由当前PC给出 \\
\code{irom\_data} & input & 32 & 指令存储器返回的32位指令 \\
\code{perip\_addr} & output & 32 & DRAM或外设访问地址，由MEM阶段ALU结果给出 \\
\code{perip\_wen} & output & 1 & 写使能信号，store指令有效 \\
\code{perip\_mask} & output & 2 & 访存宽度控制，00表示1字节，01表示2字节，10表示4字节 \\
\code{perip\_wdata} & output & 32 & 写入DRAM或外设的数据 \\
\code{perip\_rdata} & input & 32 & DRAM或外设返回的数据 \\
\bottomrule
\end{tabularx}
\end{table}

\subsection{数据通路各模块的设计}

\subsubsection{PC模块}

\code{PC}模块用于保存当前取指地址。复位时，正常上板模式下PC初始化为\code{0x8000\_0000}，以匹配企业模板中IROM起始地址；开启\code{ENABLE\_DEBUG\_TRACE}调试宏时，复位值为\code{0x0000\_0000}，便于仿真跟踪。正常运行时，PC在时钟上升沿更新为\code{NPC}模块生成的下一条指令地址。

\subsubsection{NPC模块}

\code{NPC}模块根据\code{npc\_op}选择下一条PC地址。\code{npc\_op=00}表示顺序执行或stall保持，\code{npc\_op=01}表示条件分支，\code{npc\_op=10}表示\code{jalr}，\code{npc\_op=11}表示\code{jal}。对于\code{jalr}，模块将ALU计算结果最低位置零，以符合RISC-V对\code{jalr}目标地址的要求。

\subsubsection{IF/ID流水线寄存器}

\code{IF\_ID}模块连接IF与ID阶段，保存取指阶段得到的PC、PC+4和指令。当发生\code{stall}时保持原值，当发生\code{flush}时清空寄存器内容，用于处理load-use冒险和控制冒险。

\subsubsection{Control控制模块}

\code{Control}模块根据指令\code{opcode}和\code{funct3/funct7}字段产生控制信号。当前控制器可识别R型、I型、Load、Store、Branch、U型和J型指令，输出包括\code{NpcOp}、\code{RegWrite}、\code{MemToReg\_sel}、\code{MemWrite}、\code{ALUSrcA}、\code{ALUSrcB}、\code{ALUControl}和\code{mask}等控制信号。

\subsubsection{IMMGEN立即数生成模块}

\code{IMMGEN}模块根据指令格式生成立即数。支持I型、S型、B型、U型和J型立即数，并按照RISC-V编码规则进行符号扩展和位拼接。其中B型和J型立即数最低位补0，以适配分支和跳转指令的目标地址计算。

\subsubsection{RegFile寄存器堆模块}

\code{RegFile}模块包含32个32位通用寄存器，支持两个读端口和一个写端口。写回阶段当\code{wen}有效且\code{waddr}不为0时写入寄存器。模块中还实现了同周期读写同一寄存器时的内部旁路逻辑，若读地址与写地址相同，则读端口优先返回当前写入数据。

\subsubsection{ID/EX流水线寄存器}

\code{ID\_EX}模块连接ID与EX阶段，传递PC、PC+4、源寄存器值、立即数、源/目的寄存器编号以及执行、访存和写回相关控制信号。当发生\code{id\_stall}或\code{id\_flush}时，模块将写使能、访存写使能、NPC控制等关键控制信号清零，从而在EX阶段插入无效操作。

\subsubsection{ALU\_src操作数选择模块}

\code{ALU\_src}模块用于选择ALU输入操作数。若\code{alu\_src}为1，则选择立即数或PC；若为0，则根据前递选择信号在寄存器原始值、MEM阶段结果和WB阶段写回数据中选择正确操作数。

\subsubsection{ALU模块}

\code{ALU}模块使用14位\code{ALUControl}信号控制具体运算。支持加法、减法、按位与、按位或、按位异或、逻辑左移、逻辑右移、算术右移、有符号比较、无符号比较以及分支条件判断。分支判断结果通过\code{isTrue}输出给冒险检测单元和NPC模块。

\subsubsection{EX/MEM流水线寄存器}

\code{EX\_MEM}模块连接EX与MEM阶段，保存ALU结果、store写数据、立即数、PC+4、目的寄存器编号和访存/写回控制信号。进入MEM阶段后，ALU结果作为\code{perip\_addr}输出，store写数据作为\code{perip\_wdata}输出。

\subsubsection{ex\_mux与mem\_mux模块}

\code{ex\_mux}用于在MEM阶段预先选择写回候选结果，包括PC+4、ALU结果和立即数。\code{mem\_mux}用于在WB阶段最终选择写回数据，当\code{regwrmux=10}时选择经过\code{Mask}处理的load数据，否则选择\code{ex\_mux}输出。

\subsubsection{Mask模块}

\code{Mask}模块用于处理load指令读出的数据。对于\code{lb}和\code{lh}执行符号扩展，对于\code{lbu}和\code{lhu}执行零扩展，对于\code{lw}直接保留32位数据。

\subsubsection{MEM/WB流水线寄存器}

\code{MEM\_WB}模块连接MEM与WB阶段，保存访存读数据、MEM阶段写回候选结果、目的寄存器编号以及写回控制信号。WB阶段根据选择结果将数据写回寄存器堆。

\subsubsection{forwarding\_unit前递单元}

\code{forwarding\_unit}比较EX阶段源寄存器\code{rs1/rs2}与MEM、WB阶段目的寄存器\code{rd}。若检测到数据相关且后续阶段写使能有效，则生成\code{2'b10}或\code{2'b11}前递选择信号，使ALU输入直接使用最新结果，减少流水线暂停。

\subsubsection{hazard\_detection\_unit冒险检测单元}

\code{hazard\_detection\_unit}主要处理两类冒险。第一类是控制冒险，当分支成立或发生\code{jal/jalr}跳转时，产生\code{flush}信号冲刷错误路径指令。第二类是load-use冒险，当EX阶段为load指令且其目的寄存器被ID阶段指令立即使用时，产生\code{stall}信号暂停流水线一个周期。

\subsection{数据通路信号表}

\begin{table}[H]
\centering
\caption{不同指令类型的数据通路信号表}
\label{tab:datapath_signals}
\scriptsize
\setlength{\tabcolsep}{3pt}
\renewcommand{\arraystretch}{1.15}
\begin{tabularx}{\textwidth}{
>{\centering\arraybackslash}p{1.35cm}
>{\raggedright\arraybackslash}p{2.4cm}
>{\centering\arraybackslash}p{0.9cm}
>{\centering\arraybackslash}p{1.1cm}
>{\centering\arraybackslash}p{1.1cm}
>{\centering\arraybackslash}p{1.35cm}
>{\centering\arraybackslash}p{1.2cm}
X}
\toprule
指令类型 & 指令示例 & Imm & ALU A & ALU B & ALU功能 & 写回 & 简要说明 \\
\midrule

R型运算 &
\texttt{add, sub, sll, slt, sltu, xor, srl, sra, or, and} &
无 & rs1 & rs2 & funct决定 & ALU &
寄存器间运算，结果写回rd。 \\

I型运算 &
\texttt{addi, slti, sltiu, xori, ori, andi, slli, srli, srai} &
I型 & rs1 & imm & funct决定 & ALU &
寄存器与立即数运算，结果写回rd。 \\

Load &
\texttt{lb, lh, lw, lbu, lhu} &
I型 & rs1 & imm & 地址计算 & MEM &
计算地址，读取DRAM/外设，经Mask扩展后写回rd。 \\

Store &
\texttt{sb, sh, sw} &
S型 & rs1 & imm & 地址计算 & 无 &
计算地址，将rs2写入DRAM/外设。 \\

Branch &
\texttt{beq, bne, blt, bge, bltu, bgeu} &
B型 & rs1 & rs2 & 条件比较 & 无 &
条件成立时跳转到PC+imm，否则顺序执行。 \\

\texttt{lui} &
\texttt{lui} &
U型 & 无 & imm & 无 & IMM &
立即数写回rd。 \\

\texttt{auipc} &
\texttt{auipc} &
U型 & PC & imm & 加法 & ALU &
计算PC+imm并写回rd。 \\

\texttt{jal} &
\texttt{jal} &
J型 & 无 & 无 & 无 & PC+4 &
PC跳转到PC+imm，PC+4写回rd。 \\

\texttt{jalr} &
\texttt{jalr} &
I型 & rs1 & imm & 地址计算 & PC+4 &
PC跳转到rs1+imm，PC+4写回rd。 \\

\bottomrule
\end{tabularx}
\end{table}

\subsection{控制器设计}

\subsubsection{控制信号表}

\begin{table}[H]
\centering
\caption{主要控制信号含义}
\begin{tabularx}{\textwidth}{p{3cm}p{3cm}X}
\toprule
控制信号 & 取值 & 含义 \\
\midrule
\code{NpcOp} & 00/01/10/11 & 顺序执行/条件分支/\code{jalr}/\code{jal} \\
\code{RegWrite} & 0/1 & 是否写回寄存器堆 \\
\code{MemToReg\_sel} & 00/01/10/11 & PC+4/ALU结果/访存数据/立即数 \\
\code{MemWrite} & 0/1 & 是否向DRAM或外设写入数据 \\
\code{ALUSrcA} & 0/1 & ALU输入A选择rs1或PC \\
\code{ALUSrcB} & 0/1 & ALU输入B选择rs2或立即数 \\
\code{ALUControl} & 14位独热/组合信号 & 控制ALU执行加、减、逻辑、移位或比较操作 \\
\code{mask} & funct3 & 指示访存宽度和load扩展方式 \\
\bottomrule
\end{tabularx}
\end{table}

\begin{table}[H]
\centering
\caption{各类指令控制信号表}
\scriptsize
\begin{tabularx}{\textwidth}{p{2cm}p{1.2cm}p{1.4cm}p{1.6cm}p{1.4cm}p{1.4cm}p{1.4cm}X}
\toprule
指令类型 & NpcOp & RegWrite & MemToReg & MemWrite & ALUSrcA & ALUSrcB & ALUControl \\
\midrule
R型 & 00 & 1 & 01 & 0 & 0 & 0 & funct决定 \\
I型运算 & 00 & 1 & 01 & 0 & 0 & 1 & funct决定 \\
Load & 00 & 1 & 10 & 0 & 0 & 1 & add \\
Store & 00 & 0 & -- & 1 & 0 & 1 & add \\
Branch & 01 & 0 & -- & 0 & 0 & 0 & funct决定比较 \\
\code{lui} & 00 & 1 & 11 & 0 & -- & -- & -- \\
\code{auipc} & 00 & 1 & 01 & 0 & 1 & 1 & add \\
\code{jal} & 11 & 1 & 00 & 0 & -- & -- & -- \\
\code{jalr} & 10 & 1 & 00 & 0 & 0 & 1 & add \\
\bottomrule
\end{tabularx}
\end{table}

\subsubsection{控制器设计说明}

控制器首先从指令中取出\code{opcode}，判断当前指令属于R型、I型、Load、Store、Branch、U型或J型。随后根据\code{opcode}与\code{funct}字段生成控制信号。对于R型和I型运算指令，\code{ALUControl}由\code{funct3}与\code{funct7[5]}共同决定；对于load/store和\code{jalr}指令，ALU执行加法用于地址计算；对于branch指令，ALU执行比较并输出\code{isTrue}；对于\code{lui}，写回数据选择立即数；对于\code{jal/jalr}，写回数据选择PC+4。

\section{RV32I CPU性能优化设计}

\subsection{五级流水线优化}

本CPU将指令执行划分为IF、ID、EX、MEM、WB五个阶段。流水线寄存器将各阶段结果保存并传递，使多条指令能够重叠执行。理想情况下，在流水线填满后每个周期可完成一条指令的写回，从而提高指令吞吐率。

\subsection{数据前递优化}

本设计通过\code{forwarding\_unit}解决大多数RAW数据冒险。当EX阶段源寄存器与MEM或WB阶段目的寄存器相同，且对应写使能有效时，前递单元直接将MEM或WB阶段结果送入ALU输入端，避免等待寄存器写回后再读取，从而减少流水线停顿。

\subsection{load-use暂停机制}

load指令的数据在MEM阶段后才能获得，如果紧随其后的指令立即使用该数据，则单纯前递无法满足时序要求。本设计通过\code{hazard\_detection\_unit}检测此类情况，并暂停PC和IF/ID寄存器，同时向ID/EX插入无效控制信号，使流水线等待一个周期后继续执行。

\subsection{控制冒险冲刷机制}

当条件分支成立或执行\code{jal/jalr}时，已取出的顺序路径指令可能无效。本设计在EX阶段得到分支/跳转结果后，通过\code{flush}清空IF/ID和ID/EX相关控制信息，防止错误路径指令继续执行。

\subsection{ALU组合逻辑优化}

ALU中将加减、逻辑、移位和比较运算分别计算，再通过\code{case(1'b1)}结构选择结果。该结构清晰地划分了不同运算路径，便于综合工具优化组合逻辑，也便于后续进一步针对关键路径进行时序分析。

\section{附录}

\subsection{代码清单}

\begin{longtable}{p{5cm}p{9cm}}
\caption{主要工程文件说明}\\
\toprule
文件名称 & 功能说明 \\
\midrule
\endfirsthead
\toprule
文件名称 & 功能说明 \\
\midrule
\endhead
\code{myCPU.sv} & CPU顶层模块，连接五级流水线各模块和外部接口 \\
\code{defines.sv} & 定义RV32I主要opcode宏 \\
\code{PC.sv} & 程序计数器模块 \\
\code{NPC.sv} & 下一PC生成模块 \\
\code{IF\_ID.sv} & IF/ID流水线寄存器 \\
\code{ID\_EX.sv} & ID/EX流水线寄存器 \\
\code{EX\_MEM.sv} & EX/MEM流水线寄存器 \\
\code{MEM\_WB.sv} & MEM/WB流水线寄存器 \\
\code{Control.sv} & 控制器模块 \\
\code{IMMGEN.sv} & 立即数生成模块 \\
\code{RegFile.sv} & 寄存器堆模块 \\
\code{ALU.sv} & 算术逻辑单元 \\
\code{ALU\_src.sv} & ALU操作数选择模块，包含前递选择 \\
\code{forwarding\_unit.sv} & 数据前递单元 \\
\code{hazard\_detection\_unit.sv} & 冒险检测单元 \\
\code{Mask.sv} & load数据扩展模块 \\
\code{ex\_mux.sv} & MEM阶段写回候选数据选择模块 \\
\code{mem\_mux.sv} & WB阶段最终写回数据选择模块 \\
\bottomrule
\end{longtable}

\subsection{工程目录结构}

\begin{verbatim}
project_root/
  sources/
    myCPU.sv
    defines.sv
    PC.sv
    NPC.sv
    IF_ID.sv
    ID_EX.sv
    EX_MEM.sv
    MEM_WB.sv
    Control.sv
    IMMGEN.sv
    RegFile.sv
    ALU.sv
    ALU_src.sv
    forwarding_unit.sv
    hazard_detection_unit.sv
    Mask.sv
    ex_mux.sv
    mem_mux.sv

\end{verbatim}

\subsection{参考资料}

\begin{enumerate}
    \item RISC-V Unprivileged ISA Specification, Version 2024.04.
    \item 戴维·A.帕特森，约翰·L.亨尼斯. 计算机组成与设计：硬件/软件接口 RISC-V 版[M]. 易江芳，刘先华，等译. 第2版. 北京：机械工业出版社，2023.

\end{enumerate}

\end{document}
